% Part: normal-modal-logic
% Chapter: syntax-and-semantics
% Section: introduction

\documentclass[../../../include/open-logic-section]{subfiles}

\begin{document}

\olfileid{nml}{syn}{int}

\olsection{مقدمة}

يتناول المنطق الموجهي \emph{القضايا الموجهية} وعلاقات اللزوم بينها.
وفيما يأتي أمثلة للقضايا الموجهية:
\begin{enumerate}
\item من الضروري أن~$2+2=4$.
\item من الممكن بالضرورة أن تمطر غدًا.
\item إذا كان من الممكن بالضرورة أن~$!A$، فمن الممكن أن~$!A$.
\end{enumerate}
وليست الإمكانية والضرورة الجهتين الوحيدتين؛ إذ تصنف روابط أحادية أخرى
أيضًا جهات، مثل ``ينبغي أن يكون~$!A$''، و``سيكون~$!A$''، و``تعرف دانا
أن~$!A$''، و``تعتقد دانا أن~$!A$''.

وأول ظهور للمنطق الموجهي في العرض التاريخي هنا عند أرسطو في
\emph{العبارة}. فقد تنبّه إلى أن الضرورة تقتضي الإمكانية دون العكس،
وإلى إمكان تعريف كل منهما بالأخرى. ومن ملاحظاته أن إمكان صدق
$!A \land !B$ يقتضي إمكان صدق~$!A$ وإمكان صدق~$!B$، ولا يلزم
العكس؛ وأنه متى كان~$!A \to !B$ ضروريًا، لزمت من ضرورة~$!A$
ضرورة~$!B$.

وأما البحث الحديث فابتدأ بعمل ك.~إ. لويس، وانتهى إلى كتاب لويس
ولانغفورد \emph{المنطق الرمزي} (1932). وكان منشأ بحثهما الاعتراض على
تمثيل الاستلزام بالشرط المادي؛ فليس~$!A \lif !B$ في رأيهما وافيًا
بمعنى «$!A$ تستلزم~$!B$». فجعلا الاستلزام بدل ذلك: «بالضرورة،
إذا كانت~$!A$ فإن~$!B$»، واختارا له الرمز~$!A \strictif !B$.
ثم ميّز لويس، وهو يفصل الخصائص بعضها من بعض، خمسة أنساق موجهية:
$\Log{S1}$، و\ldots، و$\Log{S4}$، و$\Log{S5}$؛ وبقي الأخيران
منها جاريين في الاستعمال.

وكان طريق لويس ولانغفورد \emph{تركيبيًا} محضًا: يضعان بديهيات
وقواعد تبدو معقولة، ثم ينظران فيما يُبرهن بها. وأما الطريق الدلالي
فلم يتيسر زمنًا طويلًا. ثم حاول رودلف كارناب في \emph{المعنى والضرورة}
(1947) أن يبنيه على \emph{وصف الحالة}، وهو مجموعة من الجمل الذرية
تُعدّ هي الصادقة في تلك الحالة. وبعد أن عمّم تعريف الصدق على الجمل
الاعتباطية~$!A$، جعل~$!A$ \emph{صادقة بالضرورة} إذا صدقت في كل
وصف للحالة. غير أن هذا البناء لا يميّز أثر الجهات \emph{المكررة}:
فالجملة من نحو «من الممكن بالضرورة \ldots من الممكن أن~$!A$»
تؤول فيه دائمًا إلى الجهة الأعمق.

ثم كان التحول الكبير بمقالة سول كريبكي «مبرهنة اكتمال في المنطق
الموجهي» (JSL 1959). وأصل الفكرة عند لايبنتس: الضروري ما يصدق
«في جميع العوالم الممكنة». لكن الإطلاق في جميع العوالم يعيد إشكال
كارناب؛ إذ إن صدق عبارة الضرورة في عالم~$w$، أو في وصف حالة~$s$،
لا يتغير بتغير~$w$. فجعل كريبكي بين العوالم علاقة \emph{وصول}~$R$.
وعندئذ تكون عبارة «بالضرورة~$!A$» صادقة في العالم~$w$ متى وفقط
متى صدقت~$!A$ في جميع العوالم~$w'$ التي \emph{يمكن الوصول إليها
من}~$w$. وكل دلالة تأخذ بصورة من هذا البناء تسمى دلالة كريبكي.
وبه اتسع البحث اتساعًا عظيمًا، فصار الكلام على منطقات موجهية متعددة.

وإذا فُسّر المنطق الموجهي بدلالة كريبكي، أمكن أن يُنظر به إلى
\emph{البنى العلائقية} «من الداخل». والمراد هنا بالبنية العلائقية
مجموعة وعلاقة ثنائية عليها؛ فمن أمثلتها طلاب الصف إذا رُتبوا بأرقام
ضمانهم الاجتماعي. ولا تنحصر هذه البنى في مجال بعينه: فقد تكون حالات
العالم مع إمكان بعضها بالقياس إلى بعض، أو حالات معرفة عامل تربطها
إمكانية معرفية، أو حالات نسق ديناميكي تربطها انتقالاته. وهذه جميعًا
تقبل النمذجة الموجهية؛ فالحالة الأولى تعطي المنطق الموجهي المعتاد،
أعني منطق جهات الصدق، وتعطي الحالات الأخرى المنطق المعرفي والمنطق
الديناميكي وما يجري مجراهما.

ونخصّ من هذا البحث بابًا يسمى «نظرية المطابقة». فمن نتائج كريبكي
المبكرة المهمة أن كثيرًا من خصائص علاقة الوصول~$R$، كالتعدي
والتناظر، يمكن وصفها \emph{في اللغة الموجهية نفسها} بمخططات مناسبة.
ومن ذلك قولهم إن انعكاسية~$R$ تطابق المخطط: «إذا كان~$!A$
ضروريًا، فإن~$!A$». وسيكون معظم نظرنا في مطابقة أنساق كلاسيكية،
مثل~$\Log{S4}$ و$\Log{S5}$، تنشأ من ضم بعض المخططات
$\Ax{D}$، و$\Ax{T}$، و$\Ax{B}$، و$\Ax{4}$، و~$\Ax{5}$ إلى بعض.

\end{document}
